% Shared preamble for the English translation of
%   P. Gabriel, "Des categories abeliennes", Bull. Soc. Math. France 90 (1962).
%
% Engine: pdflatex (also builds under xelatex/lualatex).
% twoside so that the running heads can alternate the way the original does;
% openany so that chapters do not force blank verso pages.
\documentclass[11pt,a4paper,twoside,openany]{book}

% ------------------------------------------------------------------ packages
\usepackage[T1]{fontenc}
\usepackage[utf8]{inputenc}
\usepackage{lmodern}
\usepackage[english]{babel}
\usepackage{amsmath}
\usepackage{amssymb}
\usepackage{mathtools}
\usepackage{amsthm}
\usepackage{enumitem}
\usepackage{fancyhdr}
\usepackage{microtype}
\usepackage{titlesec}
\usepackage{tikz-cd}
\usepackage[a4paper,margin=2.6cm]{geometry}
\usepackage{xcolor}
\usepackage{hyperref}

\hypersetup{
  colorlinks=true,
  linkcolor=black!70!blue,
  citecolor=black!70!blue,
  urlcolor=black!70!blue,
  pdftitle={Abelian Categories},
  pdfauthor={Pierre Gabriel},
  pdfsubject={English translation of Des categories abeliennes, Bull. Soc. Math. France 90 (1962), 323--448},
}

% ------------------------------------------------------------------ numbering
% Chapters are numbered with Roman numerals (I--VI), as in the original.
\renewcommand{\thechapter}{\Roman{chapter}.}
% Sections carry a trailing point, as in the original table of contents.
\renewcommand{\thesection}{\arabic{section}.}

% ------------------------------------------------------------------ theorems
% Gabriel numbers every result (theorem, lemma, proposition, corollary)
% with a single running counter; we reproduce that.
\newcounter{thm}
\newcounter{defn}
\newtheorem{theorem}{Theorem}[thm]
\newtheorem{lemma}[thm]{Lemma}
\newtheorem{proposition}[thm]{Proposition}
\newtheorem{corollary}[thm]{Corollary}
\theoremstyle{definition}
\newtheorem{definition}[defn]{Definition}
\newtheorem{remark}[defn]{Remark}

% Note: these counters deliberately do NOT reset at each chapter.  Gabriel
% numbers every theorem, lemma, proposition and corollary with a single
% running counter across the whole article, and the translation follows him.

% ------------------------------------------------------------------ headings
% The original prints section headings in bold, numbered, run in with a dash
% into the first paragraph.  emulating that keeps the running text close to
% the original, which matters when reading the two side by side.
\titleformat{\section}{\normalfont\normalsize\bfseries}{\thesection}{0.6em}{}
\titlespacing*{\section}{0pt}{1.2\baselineskip}{0.6\baselineskip}

% ------------------------------------------------------------------ headers
% Mirrors the original: even (verso) pages carry the author, odd (recto)
% pages carry the article title.
\pagestyle{fancy}
\fancyhf{}
\setlength{\headheight}{14pt}
\fancyhead[LE]{\thepage}
\fancyhead[RE]{\MakeUppercase{\small P. Gabriel}}
\fancyhead[LO]{\MakeUppercase{\small Abelian Categories}}
\fancyhead[RO]{\thepage}
\renewcommand{\headrulewidth}{0pt}
\fancypagestyle{plain}{\fancyhf{}\fancyhead[LO]{%
  \MakeUppercase{\small Abelian Categories}}\fancyhead[RE]{\thepage}%
  \renewcommand{\headrulewidth}{0pt}}

% ------------------------------------------------------------------ lists
\setlist[itemize]{leftmargin=1.4em,itemsep=0.25em,topsep=0.3em}
\setlist[description]{leftmargin=1.6em,labelwidth=1.3em,labelsep=0.3em,
  font=\normalfont,itemsep=0.4em,topsep=0.4em}

% ------------------------------------------------------------------ notation
% Grothendieck universes and the power-set operator.
\newcommand{\Univ}{\mathfrak{U}}
\newcommand{\Powerset}{\mathfrak{P}}
% Objects, categories and the abbreviations Gabriel fixes in Chapter I.
\newcommand{\Ob}{\mathbf{Ob}}
\newcommand{\Cat}{\mathbf{C}}
\newcommand{\Ens}{\mathbf{Ens}}
\newcommand{\Ab}{\mathbf{Ab}}
% Natural transformation, written with a dot in the original.
\newcommand{\nat}{\mathinner{\centerdot}}
\newcommand{\natt}[1]{\dot{#1}}

% ------------------------------------------------------------------ editorial
% Placeholder printed in place of a passage that has not been translated yet.
\newcommand{\notyet}[1]{%
  \par\noindent
  \begingroup
  \setlength{\fboxsep}{6pt}%
  \fcolorbox{black!40}{black!4}{%
    \parbox{\dimexpr\linewidth-2\fboxsep-2\fboxrule\relax}{%
      \small\itshape
      Not yet translated --- #1\par}}%
  \endgroup\par}